\paragraph{一般连续自同态的有限覆盖度与覆盖半群的算术谱}
在一般连续自同态的核与余核已经由 \(S\)-外分子完全刻画之后，还可把其拓扑非可逆性进一步压缩为覆盖度层面的两条闭式结论。

\leanverified{paper\_xi\_localized\_solenoid\_nonzero\_endomorphism\_finite\_cover\_degree}
\begin{corollary}[非零连续自同态必为有限覆盖，且覆盖度等于 \texorpdfstring{$S$}{S}-外分子]\label{cor:xi-localized-solenoid-nonzero-endomorphism-finite-cover-degree}
沿用定理 \ref{thm:xi-localized-solenoid-endomorphism-kernel-cokernel-sfree-numerator} 的记号。则 \(\alpha_x\) 是满射有限覆盖，并且
\[
\boxed{\;
\ker(\alpha_x)\cong \ZZ/v\ZZ,
\qquad
\deg(\alpha_x)=v.\;}
\]
特别地，\(\alpha_x\) 为同构当且仅当 \(v=1\)，亦即 \(x\in G_S^\times\)。
\end{corollary}
\begin{proof}
定理 \ref{thm:xi-localized-solenoid-endomorphism-kernel-cokernel-sfree-numerator} 已给出
\[
\ker(\alpha_x)\cong \ZZ/v\ZZ
\]
以及 \(\alpha_x\) 为同构当且仅当 \(v=1\)。
另一方面，由 \(x\neq 0\) 与 \(G_S\subset \QQ\) 无挠，离散端乘法映射
\[
m_x:G_S\to G_S
\]
是单射，故其对偶 \(\alpha_x\) 为满射。

再证有限覆盖性。由于 \(\ker(\alpha_x)\) 有限，可取单位元的开邻域 \(U\subset \Sigma_S\) 满足
\[
(U-U)\cap \ker(\alpha_x)=\{0\}.
\]
于是对不同的 \(\eta,\eta'\in \ker(\alpha_x)\)，平移集 \(U+\eta\) 与 \(U+\eta'\) 两两不交。又因 \(\alpha_x\) 为群同态，
\[
\alpha_x(U+\eta)=\alpha_x(U)
\qquad
(\eta\in \ker(\alpha_x)),
\]
并且若 \(u_1,u_2\in U\) 满足 \(\alpha_x(u_1)=\alpha_x(u_2)\)，则
\[
u_1-u_2\in (U-U)\cap \ker(\alpha_x)=\{0\},
\]
故 \(\alpha_x|_U\) 单射。由于连续满射群同态是开映射，\(\alpha_x(U)\) 为开集，且
\[
\alpha_x|_U:U\longrightarrow \alpha_x(U)
\]
为同胚。再利用平移共轭，可知每个 \(\alpha_x|_{U+\eta}\) 都是到 \(\alpha_x(U)\) 的同胚，故在该开邻域上给出恰好 \(|\ker(\alpha_x)|\) 个局部分支。于是 \(\alpha_x\) 是度数为
\[
\card{\ker(\alpha_x)}=v
\]
的有限覆盖。证毕。
\end{proof}

\begin{corollary}[非零连续自同态覆盖度半群的严格算术分类]\label{cor:xi-localized-solenoid-nonzero-endomorphism-degree-semigroup}
\leanverified{paper\_xi\_localized\_solenoid\_nonzero\_endomorphism\_degree\_semigroup}
沿用定理 \ref{thm:xi-localized-solenoid-endomorphism-kernel-cokernel-sfree-numerator} 的记号。定义
\[
\deg:\operatorname{End}_{\mathrm{cts}}(\Sigma_S)\setminus\{0\}\to \NN_{\ge 1},
\qquad
\deg(\alpha):=\card{\ker(\alpha)}.
\]
则其像恰为
\[
\boxed{\;
\operatorname{Im}(\deg)
=
\left\{
n\in \NN_{\ge 1}:\ \gcd\!\left(n,\prod_{p\in S}p\right)=1
\right\}.\;}
\]
并且对任意非零连续自同态 \(\alpha_x,\alpha_y\) 都有
\[
\boxed{\;
\deg(\alpha_x\circ \alpha_y)=\deg(\alpha_x)\deg(\alpha_y).\;}
\]
因此，\(\Sigma_S\) 的全部非平凡连续覆盖半群与去掉 \(S\)-素数后的正整数乘法半群完全同构。
\end{corollary}
\begin{proof}
由定理 \ref{thm:xi-localized-endomorphism-ring-solenoid-indecomposable}，
\[
\operatorname{End}_{\mathrm{cts}}(\Sigma_S)\cong G_S,
\]
故每个非零连续自同态都唯一写成某个 \(\alpha_x\)。

任取
\[
x=\frac{a}{b}\in G_S\setminus\{0\},
\qquad
\gcd(a,b)=1,
\qquad
b\in \langle S\rangle.
\]
把
\[
|a|=uv
\]
写成 \(u\) 的全部素因子都属于 \(S\)、而 \(v\) 与 \(\prod_{p\in S}p\) 互素的分解。由推论 \ref{cor:xi-localized-solenoid-nonzero-endomorphism-finite-cover-degree}，
\[
\deg(\alpha_x)=v.
\]
因此每个覆盖度都与 \(S\) 互素。

反之，若 \(n\in\NN_{\ge1}\) 与 \(\prod_{p\in S}p\) 互素，则 \(n\in G_S\)，并且它的 \(S\)-外分子就是 \(n\) 本身。故
\[
\deg(\alpha_n)=n.
\]
从而上述集合中的每个正整数都可实现，这就给出 \(\operatorname{Im}(\deg)\) 的精确刻画。

最后证明乘法性。设
\[
x=\varepsilon_x\frac{u_xv_x}{b_x},
\qquad
y=\varepsilon_y\frac{u_yv_y}{b_y},
\]
其中 \(\varepsilon_x,\varepsilon_y\in\{\pm1\}\)，\(u_x,u_y,b_x,b_y\) 的全部素因子都属于 \(S\)，而 \(v_x,v_y\) 都与 \(\prod_{p\in S}p\) 互素。则
\[
xy=\varepsilon_x\varepsilon_y\frac{u_xu_y}{b_xb_y}\,v_xv_y.
\]
由于 \(\frac{u_xu_y}{b_xb_y}\in G_S^\times\)，故 \(xy\) 的 \(S\)-外分子恰为 \(v_xv_y\)。再由交换性，
\[
\alpha_x\circ \alpha_y=\alpha_{xy},
\]
于是
\[
\deg(\alpha_x\circ \alpha_y)
=
\deg(\alpha_{xy})
=
v_xv_y
=
\deg(\alpha_x)\deg(\alpha_y).
\]
证毕。
\end{proof}

\endinput
